% Generated draft: requires mathematical and source review.
\documentclass[11pt]{article}
\usepackage[margin=25mm]{geometry}
\usepackage{fontspec}
\setmainfont{Latin Modern Roman}
\usepackage{amsmath,amssymb,mathtools,unicode-math}
\setmathfont{Latin Modern Math}
\usepackage{enumitem,needspace,xurl,hyperref}
\setlength{\parindent}{0pt}
\setlength{\parskip}{5pt}
\setlength{\emergencystretch}{3em}
\newcommand{\N}{\mathbb{N}}
\newcommand{\Z}{\mathbb{Z}}
\newcommand{\Q}{\mathbb{Q}}
\newcommand{\R}{\mathbb{R}}
\newcommand{\C}{\mathbb{C}}
\newcommand{\F}{\mathbb{F}}
\newcommand{\A}{\mathbb{A}}
\newcommand{\PP}{\mathbb{P}}
\newcommand{\E}{\mathbb{E}}
\newcommand{\eps}{\varepsilon}
\newcommand{\dd}{\,\mathrm{d}}
\begin{document}
\section*{UnsolvedMath ID: 20000570}
\textbf{Source problem code:} catalogue entry
\subsection{Definitions and mathematical statement}
The source is a fragment asking whether repeated identity copies with zero padding give the worst root bounds. It does not name a polynomial, matrix constraints, normalization, dimension or root statistic. The supplied report identifies it with an unnumbered Marcus--Spielman--Srivastava largest-root conjecture for mixed characteristic polynomials; that identification is a stated interpretation.

In this interpretation take positive semidefinite Hermitian \(d\)-by-\(d\) matrices \(A_1,\ldots,A_m\) satisfying
\[
 \sum_iA_i=I_d,\qquad \operatorname{tr}A_i\le\epsilon,
\]
and define
\[
 \mu[A_1,\ldots,A_m](x)=
 \left.\prod_{i=1}^m(1-\partial_{z_i})
 \det\left(xI_d+\sum_i z_iA_i\right)\right|_{z_1=\cdots=z_m=0}.
\]
The interpreted conjecture says that capped scalar identity matrices, one residual scalar matrix if necessary, and zeros maximize the largest root among feasible families with the given parameters. It is not a coordinatewise maximization of every ordered root.

The supplied candidate theorem is restricted to \(d=2\), with every fixed \(m\ge1\), \(\epsilon>0\) satisfying \(m\epsilon\ge2\). Set
\[
 q=\lfloor2/\epsilon\rfloor,\qquad r=2-q\epsilon,\qquad t_i=\operatorname{tr}A_i.
\]
It asserts, without requiring commutativity,
\[
 \mu(x)=(x-1)^2-\sum_i\det A_i,
\]
so the largest root is \(1+\sqrt{\sum_i\det A_i}\). Its exact deficit identity is
\[
 \frac{q\epsilon^2+r^2}{4}-\sum_i\det A_i
 =\frac{q\epsilon^2+r^2-\sum_i t_i^2}{4}
 +\frac12\sum_i\left\|A_i-\frac{t_i}{2}I_2\right\|_F^2\ge0,
\]
where \(\|B\|_F^2=\operatorname{tr}(B^*B)\). The first term is the trace-packing loss and the second the anisotropy loss. The reported unique largest-root maximizer, up to permutation and zero padding to \(m\) entries, consists of \(q\) copies of \((\epsilon/2)I_2\) and, when \(r>0\), one copy of \((r/2)I_2\). Feasibility supplies enough entries for these nonzero matrices. Its largest root is \(1+\tfrac12\sqrt{q\epsilon^2+r^2}\).

The report leaves the exact extremizer problem in \(d\ge3\) open in the literature it checked. It also reports a displayed counterexample to the stronger all-roots reading, but that example is not included in this record. The dimension-two largest-root theorem must therefore remain distinct from both the higher-dimensional conjecture and an all-ordered-roots claim.
\subsection{Short English statement}
The fragment is interpreted as a conjecture that scalar identity blocks and zeros maximize the largest mixed-characteristic root. The supplied report proves that interpretation for arbitrary positive semidefinite matrices in dimension two via an exact deficit identity, while retaining higher dimensions as open and rejecting the stronger all-roots interpretation.
\subsection{Sources}
\begin{itemize}
\item Local UnsolvedMath registry, permanent ID \(20000570\), source code \texttt{catalogue entry}; AIM analysis corpus, \texttt{aim-analysis-notes.json}, notes[101]. Attribution: Ulam AI / UnsolvedMath Contributors, CC BY 4.0.
\item AIM \emph{Beyond Kadison--Singer: paving and consequences}, December 1--5, 2014, source item b29: \url{https://docs.google.com/document/d/1XRw0WyG5PshqboCp99tovCRD1GCK0jXZiNLH_nw618c/edit}. No proposer or detailed conjecture is supplied in the literal fragment; original remarks and literature fields are empty.
\item The report identifies an unnumbered Marcus--Spielman--Srivastava conjecture but supplies no full citation. Its counterexample to the coordinatewise all-roots interpretation is not reproduced here.
\item Registry dated August 14, 2026: statement exact, status \texttt{partially\_solved}, classification \texttt{PARTIAL-PROGRESS}; attempt 1 records problem status open, partial result, rigor proved, novelty confidence low. The dimension-two theorem and higher-dimensional assessment are machine-generated research aids, not peer reviewed, requiring independent verification before citation.
\end{itemize}
\end{document}
